% !TeX root = main.tex
\documentclass[main.tex]{subfiles}

\begin{document}

% Max WC = 3,482

\noindent
In this research, we assess the suitablility of deep learning techniques for static malware detection, and the resiliance of these approaches to obfuscation. We evaluate three approaches to this problem, consisting of linear machine learning models, transformer based sequence classification, and an image classification model using a Convolutional Neural Network. 

Linear models comprise our baseline against which all other approaches are compared. We test a logistic regression, decision trees, and support vector machine with a radial basis function (RBF) kernel. Each of these models are trained on information extracted from the import tables of programs in our dataset, informing the models as to the external functions and libraries used by a program. 

Our sequence classification approach uses modern natural language processing techniques, such as transformer based langauge models and pretraining. We use a BERT based architecture, trained on sequences of assembly instructions. We assess the impact of pretraining by training two sequence classification models, one using a pretrained masked language model, and another trained from a random initialisation. We also train another pretrained model which has had all parameters in its attention layers frozen to assess if the representations of the pretrained model are sufficient at identifying malicious opcode sequences. 

Convolutional Neural Networks, specifically the EfficientNet architecture, offer a more computationally efficient approach trained on image representations of programs in our dataset. We test this model, and a model trained on a subset of the dataset with a more similar distribution in the size of programs between classes. This was conducted to determine if the model uses rescaling artifacts to distinguish between malicious and benign programs, as the latter are typically larger than malware. 

We then test these models on a secondary dataset of un-obfuscated programs, and increase the presence of obfuscation in each class to see if the presence of obfuscation indicative of malware. In this chapter, we provide our interpretation of the results obtained in our investigation, its relation to the work of others, and the limitations of this research. 

\section{Interpretation of Results}

\subsubsection{Baseline}

\begin{figure}
    \centering 
    \begin{subfigure}[b]{0.45\textwidth}
        \centering 
        \includegraphics[width=\textwidth]{tsne-dll.png} 
    \end{subfigure}
    \hfill 
    \begin{subfigure}[b]{0.45\textwidth}
        \centering 
        \includegraphics[width=\textwidth]{tsne-function.png} 
    \end{subfigure}
    \caption{Comparison between the functions and DLLs used by malicious and benign programs.}
    \label{fig:baseline-tsne}
\end{figure}

Our baseline approaches achieve an average classification accuracy of 90\%. Between models trained on the usage of functions and libraries, we see an marginal increase in performance of approximately 4\%. This is seen across each of the algorithms tested, and indicates that the external functions used by a programs are a better feature than external library usage information. We believe that this is due to the difference in the amount of information encoded by their respective feature vectors. Across our entire dataset, we observe 196 unique DLLs, as opposed to 5926 external functions seen more than once. Therefore, the external function features encode significantly more information than the DLL feature vector, which allow the model to better understand malicious characteristics. We also find that DLLs are more commonly seen in both malicious and benign programs as opposed to external functions, as seen in \cref{fig:baseline-tsne}. This suggests that the external libraries used by malware are more similar too those used by benign programs, meaning they are less useful as a feature. 

Of these baseline methods, we find that the decision tree algorithm achieves the best classification performance of approximately 93\%. In comparison to \citep{schultzDataMiningMethods2001a}, in which the DLL and function usage was used to identify novel malware with a rule based classification algorithm known as RIPPER, our results demonstrate improved performance. For example, the DLL usage achieved 83\% accuracy, and 89\% for function usage which seems to support our findings, despite being lower overall. We believe their results are a result of the substantially smaller dataset, which has a negative impact on model performance. They also test RIPPER on a modified version of our features, where the DLL and function level information are combined, and the number of occurances throughout the program are counted, that achieves similar performance to the function level classifier.

We were unable to assess the resiliance of these models to obfuscation, as the tool used to produce obfuscated executables hides the import table used for feature extraction. Therefore, these models would be unable to classify UPX-obfuscated executables. 

\subsubsection{Sequence Classification}

Our sequence classification model demonstrates performance that is in line with, and at times worse than the baseline approaches in sequence level tasks. For example, the pretrained model is able to identify malicious opcode sequences approximately 90\% of the time. However, this increases to approximately 95\% when the sequence level class logits are aggregated across the program. This outperforms our baseline by a small margin, but demonstrates that this approach can accurately identify malware. 

This performance, is only seen in the pretrained model, as both the model trained from a random initialisation, and the pretrained model with frozen attention weights achieve accuracy in line with, and below the baseline respectively. In the case of the model trained from a random initialisation, this may be due to insufficient training, as the pretrained model has been trained for longer, due to the pretraining step. We see in \cref{fig:rsa-pretrained-scratch} that the representations at the initial layers are highly similar between the pretrained and non-pretrained models, which decreases with depth. It may be that with increased training time, the model approaches the solution of the pretrained model. 

When the weights of the attention layers are frozen, allowing the model to modify only the classification head parameters, we see a drastic decrease in performance. We believe this is a result of reduced model complexity, as only approximately 1\% of the models parameters are updated during training. It is well established that the number of a models learnable parameters is directly associated with its performance~\citep{kaplanScalingLawsNeural2020}. By freezing the weights in the attention layers, the number of learnable parameters are reduced, leading to worse performance. \citep{leeWhatWouldElsa2019} show that when fine-tuning a model, freezing of the first 75\% of layers in similar models incurs only a 10\% decrease in performance against its fully-trainable equivalent. Therefore, a better strategy for selecting which layers to freeze in fine-tuning tasks is necessary. 

In terms of obfuscation resiliance, we find that each of these models demonstrate a high reliance on the presence of obfuscation as an indicator of malware. This can be seen in \cref{fig:base-obf-results}, where we see accuracy increase with an increase in presence of obfuscated malicious samples. In direct comparison to this, as the presence of obfuscated benign samples increases, we see only a marginal increase in accuracy. Furthermore, we also see an increase in the confidence of the models predictions as obfuscation increases, which suggests the model is heavily reliant on obfuscation. 

The same models trained on augmented data, with both fully, and partially obfuscated benign samples, demonstrate that data augmentation alone is not sufficient to increase obfuscation resiliance. For example, the model trained on the fully obfuscated benign dataset, shown in \cref{fig:seq-obf-results-ob} demonstrates an increase in accuracy with increased presence of obfuscation in the benign samples. We believe this is because the same tool was used to generate obfuscated samples in both the obfuscated benign dataset, and the obfuscation experiment. The model has discovered that characteristics of executables packed by UPX are indicative of benign samples, and therefore as we increase the number of obfuscated benign samples, we see a corresponding increase in its accuracy. This implies that UPX alone is insufficient to produce a dataset with forms of obfuscation representative of the forms used by the malicious samples in our dataset. We chose to use UPX in both cases due to its widespread support, increasing the number of samples able to be used in this experiment, however future work could use other obfuscation tools, such as Alcatraz\footnote{\url{https://github.com/weak1337/Alcatraz}} or Lime\footnote{\url{https://github.com/NYAN-x-CAT/Lime-Crypter}}.

When we train this model on the partially obfuscated benign dataset, we see a similar trend (\Cref{fig:seq-obf-results-pob}). However, as opposed to the fully obfuscated models, we do see that accuracy increases when both exclusively benign samples, and all samples are obfuscated, but a decrease when only malicious samples are obfuscated. This further indicates that the model has learned to identify UPX packing as a potential indicator of benign executables. We also see a difference in characteristics between models trained on this dataset, which is not seen across any of the other models. The model trained from a random initialisation exibits a consistent, marginal increase in accuracy across all obfuscation targets, whereas the others show a significant increase when all samples, and exclusively benign are targets. Similar characteristics across all other sequence classification models are not seen, as the trends between models are largely similar, and vary soley in the extent. 

Furthermore, as opposed to models trained on the base dataset, the obfuscated benign, and partially obfuscated benign models demonstrate decreased confidence in their classification. We see in \cref{fig:seq-obf-certainty-base}, that as presence of obfuscated samples increases, the model becomes increasingly confident in its predictions. This suggests the model has identified it as an indicator of a programs class. The decreasing confidence in the other models however, imply that obfuscation negatively impacts the models predictions. 

\subsubsection{Image Classification}

\begin{figure}
    \centering 
    \begin{subfigure}{0.8\textwidth}
        \centering 
        \includegraphics[width=\textwidth]{avg-img-rep.png}
        \caption{Mean representation of executables across classes in the full dataset.}
    \end{subfigure}
    \begin{subfigure}{0.8\textwidth}
        \centering 
        \includegraphics[width=\textwidth]{avg-img-rep-filtered.png}
        \caption{Mean representation of executables across classes in the full dataset.}
    \end{subfigure}
    \caption{}
    \label{fig:mean-img-rep}
\end{figure}

Both CNN malware detection models achieve similar performance approximately equal to the best of our baseline approaches. Therefore, these models do not achieve any notable benefit over these conventional approaches. 

The model trained on the full dataset demonstrates a marginal increase in classification performance over the filtered dataset of approximately 0.1\%. This suggests that rescaling artifacts have little impact on the solutions discovered by these models. However, as shown in \cref{fig:mean-img-rep}, there is very little difference between the average image-representations of each class in both datasets. Furthermore, we note that in both datasets, the benign samples demonstrate groups of lighter pixels at the same location. These may be artifacts originating from rescalling, which suggest that the genetic algorithm used to identify the subset is suboptimal and that further refinement is necessary. 

We also note that the image classification approaches are significantly more vulnerable to obfuscation, as an increase in accuracy is seen as we increase the number of obfuscated malicious samples. However, as we increase the number of obfuscated benign samples, we see a decrease in the accuracy of the model. This suggests heavily that the model learns to identify obfuscation rather than other features indicative of malware. 

\section{Implications \& Limitations}

\subsubsection{Implications}

We believe that this research demonstrates the potential of natural language processing techniques for not only malware detection, but reasoning about executables as a whole. We have shown that language models are capable of identifying malicious sequences of instructions, based exclusively on opcodes. It may be that increasing the vocabulary, and scale of transformer models may be able to interpret and understand these executables in a similar manner to how large language models can understand nautral langauges. 



\subsubsection{Limitations}

We identify a number of potential limitations of this work, which may have impacted the effectiveness of our work. To ensure this research is of high quality, we acknowledge these limitations, discuss how they impact the our findings, and how they may be prevented in future work. 

The first of these issues relate to the data used to train, test, and conduct our obfuscation experiment. As explained in our literature review, there are no suitable large scale datasets containing executables for malware detection. Therefore, we curate our own dataset from executables found on the hardware of the researchers, and malware sample repositories. However, we were unable to gather a sufficiently large dataset. As shown in \cref{fig:dataset:class-dist}, we are limited to approximately 10000 samples, to ensure a balanced class distribution. It is well known that greater availability of training data is associated with greater model performance. Therefore, our work is limited by the availability of this data. 

We have considered gathering benign samples from alternative sources, such as windows package repositories. However, these services often don't provide the executable, instead distributing an installer which would require extraction and installation by hand. Other sources used in the literature involve web scraping of sites such as github, and portable freeware. However, due to ethical and legal issues, we chose not to do this. 

Furthermore, we postulate that our secondary dataset, used to conduct the investigation into obfuscation may not be sufficient for our purposes for a number of reasons. The first of these relate to the extent of obfuscation, we assume that the executables within this dataset are free of obfuscation. However, when applied to this dataset the static code coverage analysis, which was previously used to determine the extent of obfuscation in our main dataset demonstrates a greater disparity between classes as shown in \cref{fig:secondary-ds-obfs-extent}. This exceeds the disparity in the primary dataset, and either indicates that code analysis coverage may not be suitable for estimating obfuscation extent within a dataset, or that the malicious executables do in fact include obfuscation. Furthermore, our analysis of the github repositories of the malware found in this dataset demonstrate a lack of diversity. We note that a majority of these samples are simple proof of concept projects\footnote{As these projects are publicly listed on github, we are able to distribute the sources of these samples, which can be found in our \href{https://github.com/Henry-Ash-Williams/masters-thesis/blob/6178028aee48cb4a8cb49f3d55f597aefc3e912d/unobfuscated-dataset.txt}{GitHub repository}}, and that keyloggers are overrepresented in this dataset. 

\begin{figure}
    \centering 
    \includegraphics[width=0.8\textwidth]{obfuscation-dist-secondary-ds.png}
    \caption{Extent of obfuscation between classes in the secondary dataset.}
    \label{fig:secondary-ds-obfs-extent}
\end{figure}

During the obfuscation experiment, we were only able to perform a single run per experiment. This means that the results may be unrepresentative of the characteristics of the approaches tested in this research. Therefore, the obfuscation experiment should be run multiple times to ensure the presence of outliers in our results are mitigated. We were unable to conduct these repreated runs due to time limitations, as a single run of the experiment took approximately an hour per model, and obfuscation target. We acknowledge that this is a limitation in our implementation, as the obfuscation and opcode extraction occurs at each sample. To make the experiment more efficient, we suggest creating a copy of the dataset, where the opcodes of obfuscated executables are extracted. 

\subsubsection{Future Work}

Future work should consider using a larger vocabulary size, and training these models on the full disassembly as opposed to the opcode sequences. Tokenization strategies such as WordPiece, which was used in the original BERT model \citep{devlinBERTPretrainingDeep2019a}, may allow for a better representation of the disassembly. We have considered using this information, however due to compute limitations, we deemed using the full disassembly to be unfeasible. 

Additionally, the impact of obfuscation on the performance of dynamic models should also be tested. We have seen a number of obfuscation techniques that target these dynamic approaches, such as \citep{liAPIASONovelAPI2023}, which prevent system API calls from being observed by the host which can prevent the extraction of API call sequences. These features are widely used in dynamic malware detection approaches seen in the literature, so this form of obfuscation may become more prevalent as malware developers become aware of these techniques. Therefore, ensuring these models are resiliant to forms of obfuscation targeting dynamic features is of high importance. 

\section{Conclusion}

In this investigation, we explore the use of deep learning techniques for static malware detection, and the resiliance of these approaches to obfuscation. Our sequence classification model, based on the BERT architecture, identifies malicious opcode sequences, which are then be aggregated across an executables disassembly to produce the file level classification. We find that this model is able to outperform our baseline approaches by a small margin, indicating that techniques from natural langauge processing are suitable in this domain, despite using a reduced model size. We also find that other natural language processing techniques, such as pretraining with masked langauge modelling has a positive effect on the performance of these models. Furthermore, we find that this approach is reliant on obfuscation to determine a programs classification, which may lead to obfuscated benign samples being misidentified as malicious in real world usage. Augmentation of training data to include obfuscated benign samples has little impact on the performance of these models under such conditions, which indicates that the forms of obfuscation found in our malicious dataset differ from the forms used in our experiments. 

Image classification approaches, offer a more computationally efficient approach. However, their performance is in line with that of our baseline. While they are more resiliant to obfuscation than the baseline approaches, as UPX hides the import table used to extract features for our baseline models, they demonstrate an increased reliance on such features for detecting malicious samples. This suggests that in real world deployment, these models would be more likely to misidentify obfuscated benign samples as malicious. Furthermore, the expected reliance of these models on artifacts originating from rescaling due to the discrepancy in size between malicious and benign samples is seemingly not an issue in this case. 

\end{document}